\chapter*{Conclusion and perspectives}
\addcontentsline{toc}{chapter}{Conclusion and perspectives}
\markboth{Conclusion and perspectives}{}

This end-of-studies internship, carried out within Qonto's Data Platform team, aimed to harden and industrialise business-critical data pipelines. The problem posed in the introduction — how to harden pipelines with no service interruption when the real source of fragility is a shared infrastructure dependency — found a coherent answer across three projects.

The first two, migrating Notion ingestion to dlt and migrating the Tableau pipelines to JWT authentication, validated the same approach: diagnosing that the coupling to Airflow's shared Docker image was the true constraint, decoupling the job into a dedicated repository and image, then removing a brittle legacy mechanism — a chain of operators with no truncation detection in one case, a token pool and an XCom hack in the other. In both cases the decisive difficulty was not writing the new code, but \emph{achieving the cutover without downstream consumers noticing}: byte-for-byte parity with production for Notion, mastery of concurrency and of the new authentication for Tableau. The third project, the CESOP transmission chain to the Dutch Digipoort gateway, applied the same rigour to a systems-integration and compliance problem, and was being finalised at the end of the internship.

In terms of results, this work eliminated two incident classes by construction, removed all manual secret rotation, took a risky dependency out of the shared image, and appreciably reduced legacy code and the number of Airflow tasks, while giving the company a new regulatory capability.

Several \textbf{perspectives} naturally extend these projects. For Notion ingestion, adding \emph{freshness tests} and volume checks would complement the resilience already gained, and dlt's Snowflake destination — inaugurated here — could be generalised to other sources currently carried by in-house operators. For Tableau, the JWT authentication template and the workbook-onboarding tooling can be reused beyond the current perimeter. For CESOP, the next step is the test campaign (\emph{ASD} connectivity then \emph{VTS} validation), the go-live, and monitoring the first real delivery through the business-validation window that may span several weeks. Finally, the \emph{zero-downtime decoupling} pattern established during this internship is a template transferable to the platform's future migrations.

On a personal level, this internship was a full immersion in production-scale data engineering: it taught me to favour root-cause diagnosis over symptom fixing, to explicitly own engineering trade-offs, and to measure the value of my work — all reflexes I consider the main takeaway of this experience.
